\centerline{\bf \Huge 10 класс}

\problem{\textbf{1}} \textbf{Критерии суммируются:} \\
\textit{4 балла} $-$ Доказано, что оставшиеся два тангенса равны 2 и 3.\\
\textit{3 балла} $-$ Наименьший угол треугольника не превосходит 60°, 
поэтому его тангенс может равняться только 1.\\
\textit{0 баллов} $-$ Любое неверное решение или неверный ответ.\\
\textit{0 баллов} $-$ Любой ответ без объяснений.\\
\textit{0 баллов} $-$ Решение отсутствует.\\
\textit{0 баллов} $-$ При обоснованных подозрениях на списывание участником из любого источника (Интернет, учебное пособие, ключи олимпиады и т.д.) член жюри имеет право указать на такие фрагменты с помощью восклицательных знаков на полях. \\

\textbf{ИЛИ:}
\textit{7 баллов} $-$ Другое полное верное решение.\\

\problem{\textbf{2}} \textbf{Критерии суммируются:} \\
\textit{0-4 балла} $-$ Решение А.\\
\textit{0-1 балл} $-$ Составление текстового условия Б.\\
\textit{0-2 балла} $-$ Любое правильное решение Б.\\
\textit{0 баллов} $-$ Любой ответ без объяснений.\\
\textit{0 баллов} $-$ Решение отсутствует.\\
\textit{0 баллов} $-$ При обоснованных подозрениях на списывание участником из любого источника (Интернет, учебное пособие, ключи олимпиады и т.д.) член жюри имеет право указать на такие фрагменты с помощью восклицательных знаков на полях. \\

\newpage

\problem{\textbf{3}} \textbf{Критерии суммируются:} \\
\textit{7 баллов} $-$ Полное верное решение.\\
\textit{4 балла} $-$ Верное решение, но без учёта возможного существования нулевых чисел.\\
\textit{0 баллов} $-$ Любой ответ без объяснений.\\
\textit{0 баллов} $-$ Решение неверное, продвижения отсутствуют.\\
\textit{0 баллов} $-$ Решение отсутствует.\\
\textit{0 баллов} $-$ При обоснованных подозрениях на списывание участником из любого источника (Интернет, учебное пособие, ключи олимпиады и т.д.) член жюри имеет право указать на такие фрагменты с помощью восклицательных знаков на полях. \\



\problem{\textbf{4}} \textbf{Критерии:} \\
\textit{7 баллов} $-$ Полное верное решение.\\
\textit{0 баллов} $-$ Любой ответ без объяснений.\\
\textit{0 баллов} $-$ Решение неверное, продвижения отсутствуют.\\
\textit{0 баллов} $-$ Решение отсутствует.\\
\textit{0 баллов} $-$ При обоснованных подозрениях на списывание участником из любого источника (Интернет, учебное пособие, ключи олимпиады и т.д.) член жюри имеет право указать на такие фрагменты с помощью восклицательных знаков на полях. \\

\problem{\textbf{5}} \textbf{Критерии:} \\
\textit{4-7 баллов} $-$ Полное верное решение.\\
\textit{1 балл} $-$ Ребро куба равно $\sqrt{6}$.\\
\textit{0 баллов} $-$ Любой ответ без объяснений.\\
\textit{0 баллов} $-$ Решение неверное, продвижения отсутствуют.\\
\textit{0 баллов} $-$ Решение отсутствует.\\
\textit{0 баллов} $-$ При обоснованных подозрениях на списывание участником из любого источника (Интернет, учебное пособие, ключи олимпиады и т.д.) член жюри имеет право указать на такие фрагменты с помощью восклицательных знаков на полях. \\